\section{Grenzen der Arbeit}

Die vorliegende Arbeit demonstriert die Machbarkeit einer kombinierten Balance- und Spurführungs-Regelung für ein 
selbstbalancierendes Fahrrad in einer Webots-Simulationsumgebung. Trotz der erzielten positiven Ergebnisse weist die entwickelte 
Lösung verschiedene Limitierungen auf, die sowohl methodischer als auch technischer Natur sind. Diese Grenzen sind teilweise 
bewusst gewählte Vereinfachungen zur Fokussierung auf die Kernproblematik, teilweise aber auch ureigene Beschränkungen 
der gewählten Ansätze. Eine transparente Darstellung dieser Limitierungen ist essentiell für die Einordnung der Ergebnisse und 
bildet die Grundlage für zukünftige Erweiterungen des Systems.

\subsection{Simulationsbedingte Limitierungen}

\subsubsection{Inertial Measurement Unit (IMU)}
Die Sensordatenerfassung in der Webots-Simulationsumgebung weist fundamentale 
Unterschiede zur realen Hardware-Implementierung auf, die erhebliche Auswirkungen auf die Übertragbarkeit der entwickelten 
Regelungsalgorithmen haben. Diese Diskrepanz manifestiert sich besonders deutlich beim Vergleich zwischen der verwendeten 
\texttt{InertialUnit} in Webots und dem in den Vorarbeiten eingesetzten BNO055-IMU-Sensor.

Während die Webots-\texttt{InertialUnit} über die API-Funktion \texttt{wb\_inertial\_unit\_get\_roll\_pitch\_yaw()} exakte 
Euler-Winkel ohne jegliche Messfehler bereitstellt, unterliegt der reale BNO055-Sensor einer Vielzahl von Fehlerquellen und 
Hardware-Limitierungen. Die Arbeiten von Ranz (2021) und Karabila (2022) dokumentieren diese Problematik ausführlich und 
zeigen die Komplexität der realen Sensordatenverarbeitung auf.

Der BNO055 arbeitet als 9-Achsen-Fusionssensor, der Accelerometer-, Gyroscope- und Magnetometer-Daten intern verarbeitet und 
über I²C-Kommunikation 16-Bit-Rohwerte ausgibt. Diese müssen über konfigurierbare Skalierungsfaktoren in physikalische Einheiten konvertiert werden \citep{ranz2021bachelor}. Die endliche Auflösung führt zwangsläufig 
zu Quantisierungsfehlern, die bei der hochfrequenten Balance-Regelung (500\,Hz) akkumulieren können. Zusätzlich erfordert der Sensor 
eine mehrstufige Kalibrierung seiner Teilkomponenten, wobei Ranz die Implementierung einer \texttt{set\_axis\_mode()}-Funktion für 
das erforderliche Achsenremapping dokumentiert.

Besonders kritisch erweist sich die Temperaturabhängigkeit des BNO055, die zu zeitvarianten Bias-Fehlern führt. Diese Drift-Effekte 
akkumulieren über die Betriebszeit und können die Regelstabilität beeinträchtigen. Im Gegensatz dazu liefert die Webots-Simulation 
zeitinvariante, temperaturunabhängige Messungen. Die I²C-Kommunikation über \texttt{/dev/i2c-2} auf dem BeagleBone Black verursacht 
zusätzlich variable Kommunikationslatenzen, die Ranz durch spezielle Timeout-Mechanismen und Pufferstrategien kompensieren musste.

Diese Diskrepanzen zwischen Simulation und Realität lassen sich durch verschiedene Ansätze verringern. Die Implementierung eines 
konfigurierbaren Rauschmodells in der Webots-Simulation könnte die spektralen Eigenschaften des BNO055 nachbilden, beispielsweise
durch Überlagerung von weißem Rauschen und charakteristischem 1/f-Rauschen\footnote{1/f-Rauschen, auch als Rosa Rauschen bezeichnet, 
ist ein frequenzabhängiges Rauschsignal, dessen Leistungsdichte proportional zu 1/f abnimmt. Es tritt charakteristisch in 
elektronischen Bauelementen und Sensoren auf und führt zu niederfrequenten Drifteffekten.}. Kommunikationslatenzen ließen sich durch zeitliche 
Pufferung der Sensordaten über mehrere Simulationsschritte simulieren. Systematische Kalibrierungsfehler und temperaturabhängige 
Drift könnten durch langsam variierende Bias-Terme modelliert werden, die von simulierten Umgebungsparametern abhängen.

%--

\subsubsection{Kamera und Vision-Pipeline}
Die Kamerasensorik zeigt analog zur IMU-Problematik erhebliche Unterschiede zwischen der 
idealisierten Webots-Simulation und realen Implementierungen. Diese Diskrepanz wird besonders deutlich bei der Betrachtung von 
Girays (2024) Vision-basierter MPC-Implementierung, die eine YOLOv8-Segmentierung für die Spurerkennung einsetzt.

Die Webots-Kamera erzeugt perfekte, rauschfreie Bilder mit konstanter Beleuchtung und ohne Bewegungsunschärfe oder optische 
Verzerrungen. Girays reale Implementierung auf dem NVIDIA Jetson Nano hingegen musste verschiedene hardware- und umgebungsbedingte 
Herausforderungen bewältigen. Das YOLOv8-Segmentierungsmodell wurde mit einem spezifischen Datensatz für vier Klassen trainiert 
(\textit{meadow}, \textit{obstacle}, \textit{street\_main}, \textit{street\_side}) und mittels TensorRT für die Inferenz optimiert, 
um Bildauswertungszyklen von etwa 100\,ms zu erreichen \citep{vision_mpc_readme}.

Während die Webots-Simulation deterministisch identische Bilder unter kontrollierten Bedingungen liefert, unterliegt die reale 
Kameraverarbeitung verschiedenen Störfaktoren. Girays Implementierung dokumentiert Herausforderungen bei der Echtzeitverarbeitung, 
wobei die Inferenzzeit des YOLOv8-Modells auf eingebetteter Hardware kritisch für die angestrebte 20\,Hz-Abtastrate der Vision-Pipeline 
wird. \pdfcomment{Hier sich sein mit den hz}Die in der Simulation problemlos erreichbare Bildverarbeitung erfordert in der Realität komplexe 
Optimierungsstrategien wie TensorRT-Beschleunigung und Modellkompression.

Die monokulare Kameraarchitektur in Girays System limitiert zusätzlich die Tiefenwahrnehmung auf geometrische Approximationen. Während 
dies in der 2D-Simulation von Webots unproblematisch ist, führt es in realen 3D-Umgebungen zu Unsicherheiten bei der Entfernungsschätzung 
und Hinderniserkennung. Girays Arbeit identifiziert diese Limitation explizit und schlägt stereo-basierte oder LiDAR-gestützte Ansätze für 
zukünftige Implementierungen vor.

Die Übertragung von Girays Vision-Pipeline in die Webots-Simulation abstrahiert von diesen realen Herausforderungen und kann daher zu 
optimistischen Performance-Bewertungen führen. Algorithmen, die in der idealisierten Simulationsumgebung robust funktionieren, können 
unter realen Bedingungen mit variierenden Lichtverhältnissen, Bewegungsunschärfe und Hardwarelimitierungen versagen.

Zusätzlich ist die IPC-Kommunikation zwischen den Controllern in Webots idealisiert und weist keine Paketverluste oder variable Latenzen 
auf. Reale Kommunikationssysteme (UART, Ethernet, CAN) können hingegen variable Latenzen je nach Systemlast, Paketverluste durch 
Übertragungsfehler sowie Jitter durch unregelmäßige Ankunftszeiten der Nachrichten aufweisen. Diese Effekte können die Stabilität 
des Gesamtsystems beeinträchtigen und erfordern robuste Kommunikationsprotokolle mit Fehlerbehandlung und Timeout-Mechanismen.

Die aufgezeigten Unterschiede zwischen idealisierter Simulation und realer Sensorik stellen eine der kritischsten Limitierungen für 
die direkte Übertragbarkeit der entwickelten Regelungsstrategien dar. Sie unterstreichen die Notwendigkeit einer mehrstufigen 
Validierungsstrategie, die von der Simulation über Hardware-in-the-Loop-Tests bis hin zu realen Fahrversuchen reicht.

\subsubsection{Vereinfachte Physikmodellierung}
Wie in Kapitel \ref{sec:ode} dargestellt, basiert die Webots-Simulation auf der Open Dynamics Engine (ODE). Die dort beschriebenen 
Eigenschaften und Implementierungsdetails von ODE führen jedoch zu verschiedenen grundlegenden Limitierungen, die Auswirkungen 
auf die Realitätstreue der Fahrradsimulation haben und die Übertragbarkeit der Ergebnisse auf reale Systeme teilweise einschränken.

Die kritischsten Limitierungen betreffen die vereinfachte Reifen-Fahrbahn-Modellierung und numerische Stabilitätsprobleme. ODE reduziert 
den Reifenkontakt auf eine Punkt-zu-Ebene-Interaktion mit konstantem Reibungskoeffizienten, während reale Reifen komplexe nichtlineare 
Charakteristika aufweisen. Besonders die fehlende Modellierung des Seitenschlupfes, der bei Kurvenfahrten die lateralen Führungskräfte 
bestimmt, beeinträchtigt die Realitätstreue der Fahrdynamik erheblich \citep{catto2009ode_friction}. Im Gegensatz zu den in der 
Fahrzeugdynamik etablierten detaillierten Reifenmodellen \citep{schramm2018fahrzeugdynamik} vernachlässigt die ODE-Simulation 
fundamentale Aspekte der Reifen-Fahrbahn-Interaktion.

Der verwendete explizite Euler-Integrator\footnote{Der Euler-Integrator ist ein numerisches Verfahren zur Lösung von Differentialgleichungen, das aufgrund seiner einfachen Implementierung weit verbreitet ist, jedoch bei großen Zeitschritten oder steifen Systemen zu Stabilitätsproblemen neigt.} kann bei der hochfrequenten Balance-Regelung (500\,Hz) zu numerischen Instabilitäten führen, 
was die Validität der Simulationsergebnisse gefährdet \citep{smith2006ode_stability}. Zusätzlich vernachlässigt die vereinfachte 
Mehrkörperdynamik-Modellierung wichtige Aspekte wie Spiel in Lagern und elastische Verformungen des Rahmens.

Diese fundamentalen Limitierungen der ODE-basierten Physikmodellierung stellen eine erhebliche Einschränkung für die Übertragbarkeit der 
Simulationsergebnisse auf reale Fahrradsysteme dar und erfordern eine kritische Bewertung der erzielten Regelungsperformance.

\textbf{Hardware-Transfer-Implikationen:}

Die Übertragung der ODE-basierten Simulationsergebnisse auf reale Hardware wird zusätzlich durch verschiedene implementierungsspezifische 
Faktoren erschwert. Die Webots-Simulation abstrahiert von vielen hardwarespezifischen Details wie PWM-Charakteristika mit ihren 
Nichtlinearitäten und Totzeiten der Motoransteuerung, mechanischem Spiel in Getrieben und Lagern sowie Temperaturdrift und Verschleiß, 
die das Langzeitverhalten und die Degradation der Komponenten beeinflussen. Diese Faktoren können das reale Systemverhalten erheblich 
von der idealisierten Simulation abweichen lassen und erfordern robuste Adaptationsstrategien bei der Hardware-Implementierung.



